\begin{pastexam}[label={px:ust-115-ee-ds-09}]{臺聯大電機類 115 學年度 第 9 題}
（單選；2 分）

Given a ``Perfect Binary Tree'' (i.e., every level is full), with a height of h (total number of
nodes n $\approx 2^{h+1}$), compare the worst-case space complexity of using a stack for depth-first
search (DFS) versus using a queue for breadth-first search (BFS).
\begin{enumerate}[label=(\Alph*)]
\item Depth-first search (DFS) is relatively memory-intensive because of its deep recursion depth.
\item BFS is relatively memory-intensive because a queue needs to store all the leaf nodes at the
lowest level.
\item Both consume roughly the same amount of memory O($n$).
\item Both consume roughly the same amount of memory O($\log n$)
\end{enumerate}
\end{pastexam}

\begin{pastsoln}
答案是 (B)。

高度 $h$ 的完美二元樹有 $n = 2^{h+1} - 1$ 個節點，$h = \Theta(\log n)$，最底層有
$2^h = (n + 1)/2$ 個葉子。

\textbf{DFS（堆疊）}：不論用遞迴還是自己的堆疊，堆疊裡放的是從根到目前節點那條路徑上的節點
（自己用堆疊、一次把兩個孩子都 push 時，再加上路徑旁邊還沒走的兄弟，每層最多一個）。 路徑最長
$h + 1$ 個節點，所以堆疊大小是 $O(h) = O(\log n)$。

\textbf{BFS（佇列）}：層序走訪（資結書第~7.3.2~節）把第 $d$ 層全部取出時，第 $d + 1$ 層的
$2^{d+1}$ 個節點都已經進了佇列。 處理完倒數第二層時，佇列裡剛好是最底層全部 $2^h = (n+1)/2$ 個
葉子，所以佇列大小是 $\Theta(n)$。

所以 BFS 比較吃記憶體，原因正是 (B) 說的要同時存最底層的葉子。
\begin{itemize}
\item (A) 完美二元樹的深度只有 $\Theta(\log n)$，DFS 的遞迴並不深。 錯。
\item (B) 對。
\item (C) DFS 只要 $O(\log n)$，不是 $O(n)$。 錯。
\item (D) BFS 要 $\Theta(n)$，不是 $O(\log n)$。 錯。
\end{itemize}
（如果樹退化成一條鏈，情況反過來：DFS 的堆疊是 $\Theta(n)$、BFS 的佇列只有 $O(1)$。 題目限定是
完美二元樹。）
\end{pastsoln}
